\slajdid{07-s017}
\begin{frame}[label=07-s017]
\gusttekst\setlength{\parskip}{4pt}
Za simulaciju dizajna: Jednom kada uđete u sistem projektovanja zasnovan na računaru, verovatno ćete želeti da simulirate rad vašeg kola da biste saznali da li će ono ispuniti funkcionalne i vremenske zahteve razvijene tokom procesa specifikacije. Ako ste kreirali jedan ili više test bench-a kao deo vaše specifikacije dizajna, onda ćete koristiti simulator da primenite test bench-a na svoj dizajn onako kako je napisan za sintezu (funkcionalna simulacija) i eventualno korišćenje verzije nakon sinteze dizajna takođe. \par\vspace{4mm}Za projektnu dokumentaciju: Karakteristike strukturiranog programiranja VHDL-a, zajedno sa njegovim karakteristikama upravljanja konfiguracijom, čine VHDL prirodnim oblikom u kojem se dokumentuje veliko i složeno kolo. Na vrednost korišćenja jezika visokog nivoa kao što je VHDL za projektnu dokumentaciju ukazuje činjenica da Ministarstvo odbrane SAD sada zahteva VHDL kao standardni format za komunikaciju zahteva dizajna između vladinih podizvođača. \par\vspace{4mm}Kao alternativa šemi: Šeme su dugo bile deo dizajna elektronskih sistema i malo je verovatno da će uskoro izumreti. Šeme imaju svoje prednosti, posebno kada se koriste za prikazivanje kola u obliku blok dijagrama. Iz tog razloga mnogi VHDL alati za dizajn sada nude mogućnost kombinovanja šematskih i VHDL reprezentacija u dizajnu.
\end{frame}
